\subsection{例：连续函数空间中的线性表示}

设 G 为在集合 X 上左作用的离散群。称 \(\chi\) 为 \(G \times X\) 上的复值函数，且若它是乘子，则

\begin{enumerate}
\def\labelenumi{(\arabic{enumi})}
\tightlist
\item
\end{enumerate}

\[
\chi (e, x) = 1 \quad \text { for   all } x \in X;\tag{2}
\]

\[
\chi (s t, x) = \chi (s, t x) \chi (t, x) \quad \text { for   all } s, t \text { in } G, x \in X.
\]

于是有

\[
\chi (t ^ {- 1}, t x) \chi (t, x) = 1 \quad \text { for   all } t \in \mathrm{G}, x \in \mathrm{X},\tag{3}
\]

特别地，\(\chi(t, x) \neq 0\) 对所有 \(t \in \mathbf{G}\)、\(x \in \mathbf{X}\) 成立。

对每个复值函数 \(f\)，定义在 \(\mathbf{X}\) 上，及每个 \(s \in \mathbf{G}\)，令 \(\pmb{\gamma}_{\chi}(s)f\) 为 \(\mathbf{X}\) 上由下式定义的复值函数：

\[
\left(\boldsymbol {\gamma} _ {\chi} (s) f\right) (x) = \chi (s ^ {- 1}, x) f (s ^ {- 1} x).\tag{4}
\]

则 \(\gamma_{\chi}(e)f = f\)，并且

\[
\begin{array}{r l} & {\big (\pmb {\gamma} _ {\chi} (s) \pmb {\gamma} _ {\chi} (s ^ {\prime}) f \big) (x) = \chi (s ^ {- 1}, x) \big (\pmb {\gamma} _ {\chi} (s ^ {\prime}) f \big) (s ^ {- 1} x)} \\ & {\qquad = \chi (s ^ {- 1}, x) \chi (s ^ {- 1}, s ^ {- 1} x) f (s ^ {- 1} s ^ {- 1} x)} \\ & {\qquad = \chi \big ((s s ^ {\prime}) ^ {- 1}, x) \big) f \big ((s s ^ {\prime}) ^ {- 1} x \big) = \big (\pmb {\gamma} _ {\chi} (s s ^ {\prime}) f \big) (x),} \end{array}
\]

因此 \(\pmb{\gamma}_{\chi}\) 是 G 的线性表示。当 \(\chi = 1\) 时，就得到自同态 \(\pmb{\gamma}(s)\)（Ch. VII, §1, No.~1, 公式 (3)）。

现在假设 G 和 X 都是局部紧的，G 在 X 上连续作用，且 \(\chi\) 在 \(G \times X\) 上连续。则 \(\mathcal{C}(X)\) 和 \(\mathcal{K}(X)\) 对 \(\boldsymbol{\gamma}_{\chi}(s)\) 稳定，从而得到 G 在 \(\mathcal{C}(X)\) 和 \(\mathcal{K}(X)\) 中的线性表示，仍记为 \(\gamma_{\chi}\)。

命题 4. --- 线性表示 \(\pmb{\gamma}_{\chi}\) 在 G 的 \(\mathcal{C}(\mathbf{X})\) 和 \(\mathcal{K}(\mathbf{X})\) 中是连续的。

映射 \((s, f) \mapsto (s, \gamma(s)f)\)，从 \(\mathbf{G} \times \mathcal{C}(\mathbf{X})\) 到 \(\mathbf{G} \times \mathcal{C}(\mathbf{X})\)，是连续的（No.~2, 引理 3）。另一方面，映射 \((s, f) \mapsto \chi(s, \cdot)f\)，从 \(\mathbf{G} \times \mathcal{C}(\mathbf{X})\) 到 \(\mathcal{C}(\mathbf{X})\)，是连续的；因为若 \(s\) 趋于 \(s_0\)（在 \(\mathbf{G}\) 中），则 \(\chi(s, \cdot)\) 趋于 \(\chi(s_0, \cdot)\)，在 \(\mathbf{X}\) 的每个紧子集上一致；此外，若 \(f\) 趋于 \(f_0\)，且二者属于 \(\mathcal{C}(\mathbf{X})\)，则 \(\chi(s, \cdot)f\) 趋于 \(\chi(s_0, \cdot)f_0\)，在 \(\mathbf{X}\) 的每个紧子集上一致，这就证明了断言。因此 \(\gamma_{\chi}\) 作为 \(\mathbf{G}\) 在 \(\mathcal{C}(\mathbf{X})\) 中的表示是连续的。

我们来证明 \(\pmb{\gamma}_{\chi}\) 作为 G 在 \(\mathcal{K}(\mathrm{X})\) 中的表示是连续的。由于 \(\mathcal{K}(\mathrm{X})\) 是 Banach 空间的归纳极限，它是桶状的（TVS, III, §4, No.~1, 命题 3 的推论 3），所以只需证明 \(\pmb{\gamma}_{\chi}\) 分别连续（No.~1, 命题 1）。现在令 H 为 X 的紧子集，令 \(s_0 \in \mathrm{G}\)。令 V 为 G 中 \(s_0\) 的紧邻域，并令 L = VH；L 在 X 中是紧的。对每个 \(f \in \mathcal{K}(\mathrm{X}, \mathrm{H})\)，\(\pmb{\gamma}_{\chi}(s_0)f\) 的支集包含于 L 中，并且

\[
\sup _ {x \in X} | (\boldsymbol {\gamma} _ {\chi} (s _ {0}) f) (x) | \leqslant \sup _ {x \in L} | \chi (s _ {0} ^ {- 1}, x) | \cdot \sup _ {x \in X} | f (x) |,
\]

因此 \(f \mapsto \boldsymbol{\gamma}_{\chi}(s_0)f\) 是从 \(\mathcal{K}(\mathrm{X},\mathrm{H})\) 到 \(\mathcal{K}(\mathrm{X},\mathrm{L})\) 的连续线性映射；由此 \(f \mapsto \boldsymbol{\gamma}_{\chi}(s_0)f\) 是从 \(\mathcal{K}(\mathrm{X})\) 到自身的连续线性映射（TVS, II, §4, No.~4, 命题 5）。另一方面，\(\mathcal{K}(\mathrm{X},\mathrm{L})\) 的拓扑由 \(\mathcal{C}(\mathrm{X})\) 的拓扑诱导。根据已经证明的结果，映射 \(s \mapsto \boldsymbol{\gamma}_{\chi}(s)f\) 从 \(\mathrm{V}\) 到 \(\mathcal{K}(\mathrm{X},\mathrm{L})\) 连续。这就完成了 \(\boldsymbol{\gamma}_{\chi}\) 分别连续的证明。

命题 5. --- 假设每个函数 \(\chi(s,\cdot)\) 都有界。则 \(\gamma_{\chi}\) 保持 \(\overline{\mathcal{K}(X)}\) 不变，并且 \(\gamma_{\chi}\) 作为 \(G\) 在 \(\overline{\mathcal{K}(X)}\) 中的线性表示是连续的。

显然 \(\pmb{\gamma}_{\chi}\) 保持 \(\overline{\mathcal{K}(\mathrm{X})}\) 不变，且每个 \(\pmb{\gamma}_{\chi}(s)\) 在 \(\overline{\mathcal{K}(\mathrm{X})}\) 中连续。另一方面，对每个 \(f \in \mathcal{K}(\mathrm{X})\)，\(s \mapsto \pmb{\gamma}_{\chi}(s)f\) 是从 \(\mathbf{G}\) 到 \(\mathcal{K}(\mathbf{X})\) 的连续映射，因而也是到 \(\overline{\mathcal{K}(\mathbf{X})}\) 的连续映射。因此 \(\pmb{\gamma}_{\chi}\) 在 \(\overline{\mathcal{K}(\mathbf{X})}\) 中的表示是连续的（No.~1, 命题 2）。

